% =============================================================================
%  Lab 1 - Design Flow and the Transparent Latch
%
%  Seven sections. Same seven every week. Do not add, remove or reorder them.
%  Delete the guidance comments as you fill each section in.
% =============================================================================

\chapter{Lab 1 --- Design Flow and the Transparent Latch}

\section{Objective}
% Two or three sentences, in your own words. What did this lab ask you to build,
% and what idea was it teaching? Do not copy the handout.

\section{Specification}
% The required behaviour as a TABLE, not as prose. Number every row, so that
% your Discussion section can refer to "row 3" instead of describing it.

\begin{table}[htbp]
\centering
\caption{Transition table (Type T) for \texttt{TransLatch}.}
\label{tab:lab01-truth}
\begin{tabular}{clll||ll}
\toprule
\# & \texttt{Clr} & \texttt{En} & \texttt{D} & $Q^{+}$ & note \\
\midrule
0 & 1 & -- & -- & \textbf{0} & asynchronous clear \\
1 & 0 & 0 & -- & \textbf{Q} & hold --- the memory row \\
2 & 0 & 1 & 0 & \textbf{0} & transparent \\
3 & 0 & 1 & 1 & \textbf{1} & transparent \\
\bottomrule
\end{tabular}
\end{table}

\subsection*{Boolean derivation}
% Three lines. Line (2) must NAME the law or K-map grouping you used.
% "It simplified" is not an answer.

\begin{align*}
\text{(1) canonical} \quad & Q^{+} = \overline{Clr}\cdot En \cdot D \;+\; \overline{Clr}\cdot\overline{En}\cdot Q \\
\text{(2) factored}  \quad & Q^{+} = \overline{Clr}\cdot\left( En\cdot D + \overline{En}\cdot Q \right)
   \qquad \text{[\emph{name the law here}]} \\
\text{(3) VHDL}      \quad & \text{see Listing~\ref{lst:code/lab01/TransLatch.vhd}}
\end{align*}

% One sentence: which term makes this memory rather than combinational logic?

\section{Design}
% The entity interface, and WHY you chose this architecture style.
% For this lab: why is a process with an if-statement the right description,
% and what does the deliberately missing 'else' mean in hardware?

\section{Implementation}
% Include the REAL .vhd file. Never paste code into this document.
% Upload your sources into code/lab01/ in this Overleaf project.

\vhdlfile{code/lab01/TransLatch.vhd}{\texttt{TransLatch} --- 1-bit transparent latch with asynchronous clear}

% \vhdlfile{code/lab01/TransLatch8.vhd}{\texttt{TransLatch8} --- 8-bit version}

\section{Verification}
% Your test vectors as a table, then the COVERAGE ARGUMENT: why is this set
% sufficient? What would a wrong design have to do to pass all of them?

\begin{table}[htbp]
\centering
\caption{Test vectors applied to \texttt{TransLatch}.}
\label{tab:lab01-vectors}
\begin{tabular}{clllll}
\toprule
\# & \texttt{Clr} & \texttt{En} & \texttt{D} & expected $Q$ & what it proves \\
\midrule
1 &   &   &   &   & \\
2 &   &   &   &   & \\
3 &   &   &   &   & \\
\bottomrule
\end{tabular}
\end{table}

% \vhdlfile{code/lab01/tb_TransLatch.vhd}{Self-checking testbench}

\section{Results}
% Evidence set for this lab: E1, E2, E7.
% EVERY figure needs one sentence saying what it PROVES. A figure with no
% sentence earns zero. Replace each \figplaceholder with \resultfig.

\figplaceholder{the RTL Viewer output}
% \resultfig{figs/lab01/rtl.png}
%   {The RTL Viewer shows a single latch primitive and no other storage,
%    confirming the design matches Table~\ref{tab:lab01-truth}.}
%   {fig:lab01-rtl}

\figplaceholder{the Questa waveform, with the hold interval annotated}
% \resultfig{figs/lab01/waveform.png}
%   {Between t = ... and t = ..., D changes while En is low and Q does not
%    follow, which is row 1 of Table~\ref{tab:lab01-truth}.}
%   {fig:lab01-wave}

% E1: state the error and warning counts, and sort every warning into
%     fix / explain / ignore.

\section{Discussion}
% REQUIRED. One bug you hit, how you found it, and what it cost you.
% "Everything worked first try" scores zero. If you genuinely had no bug,
% write about something that surprised you in the RTL Viewer or the waveform.
%
% One line linking this lab to lecture, e.g.:
% "This is the storage element behind the register discussed in Homework 1."
